% !TeX root = ../偏微分方程笔记.tex
% 由主文档通过 \input 载入；请编译主文档。
\setcounter{chapter}{2} % 保留原书章号。
\chapter{泛函分析工具}
\label{app:functional}

\section{巴拿赫空间与希尔伯特空间}
完备的赋范线性空间称\emph{巴拿赫空间（Banach space）}；内积诱导的范数使之完备时称\emph{希尔伯特空间（Hilbert space）}。完备性保证柯西近似真正收敛到所选空间中的元素。

希尔伯特空间 $H$ 中任意闭线性子空间 $M$ 都有正交投影：每个 $x\in H$ 唯一分解为 $x=m+z$，$m\in M$、$z\perp M$。证明可取距离极小化序列，利用平行四边形恒等式说明它为柯西列，再由闭性取得极小点；沿 $M$ 中任意方向求一元二次函数的极小值得正交性。

\begin{theorem}[里斯表示定理（Riesz representation theorem）]
在实希尔伯特空间上，每个连续线性泛函 $F$ 都唯一写成 $F(v)=(z,v)_H$，其中 $z\in H$，且 $\|F\|=\|z\|_H$。
\end{theorem}
\begin{proof}
$F=0$ 时取 $z=0$。否则 $\ker F$ 是闭子空间且其正交补为一维：任选 $F(x)\ne0$，每个 $v$ 减去 $F(v)x/F(x)$ 就落在核中。取其正交补的单位向量 $e$，令 $z=F(e)e$，分解 $v$ 即验证表示。范数等式来自柯西--施瓦茨和取 $v=z/\|z\|$，唯一性由差向量与所有向量正交得到。
\end{proof}

\section{对偶空间与弱收敛}
赋范空间 $X$ 的连续线性泛函构成对偶空间 $X^*$，范数为 $\|F\|_{X^*}=\sup_{\|v\|_X\le1}|F(v)|$，作用也记作 $\langle F,v\rangle$。若 $F(x_j)\to F(x)$ 对每个 $F\in X^*$ 成立，称 $x_j$ \emph{弱收敛（weak convergence）}到 $x$，记 $x_j\rightharpoonup x$。若 $F_j\in X^*$ 对每个固定 $x\in X$ 都有 $F_j(x)\to F(x)$，称为\emph{弱星收敛（weak-star convergence）}。

可分希尔伯特空间中的有界序列有弱收敛子列：对一组可数规范正交基逐个提取系数收敛子列，再对角选取；有界性使极限系数平方可和，并可用有限维逼近控制任意测试向量。注意这并不保证范数收敛。例如规范正交序列弱收敛到零，但各项范数都为一。

范数弱下半连续。在希尔伯特空间中，若 $x_j\rightharpoonup x\ne0$，则 $(x_j,x)\to\|x\|^2$，由柯西--施瓦茨得 $\|x\|\le\liminf\|x_j\|$；$x=0$ 时结论直接成立。结合范数收敛，弱收敛可升级为强收敛，因为展开 $\|x_j-x\|^2$ 即可验证。

\section{抽象工具的使用清单}
\begin{theorem}[巴拿赫不动点定理（Banach fixed-point theorem）]
设非空完备度量空间 $X$ 上的映射 $T:X\to X$ 满足 $d(Tx,Ty)\le qd(x,y)$，其中 $0\le q<1$。则 $T$ 有唯一不动点，任意初值的迭代 $x_{m+1}=Tx_m$ 都收敛到它。
\end{theorem}
\begin{proof}
连续迭代差满足 $d(x_{m+1},x_m)\le q^m d(x_1,x_0)$，等比级数控制任意尾段，故迭代为柯西列。完备性给出极限 $x_*$，利普希茨连续性给出 $Tx_*=x_*$。若另有 $y_*$，则 $d(x_*,y_*)\le qd(x_*,y_*)$，所以两者相同。
\end{proof}
在函数闭球中使用时，既要证明压缩，也要证明映射确实留在闭球内。拉克斯--米尔格拉姆定理则见第\ref{chap:elliptic-weak}章，其对象是希尔伯特空间上的有界强制双线性型，不能用压缩条件替代。

另一个常用结论是紧自伴算子的谱定理：若 $K:H\to H$ 紧且自伴，则非零谱值为实本征值，每个非零本征空间有限维，非零本征值的唯一可能聚点为零，并且 $H$ 是 $\ker K$ 与这些本征空间闭线性张成的正交直和。这里紧性指有界序列的像具有范数收敛子列，自伴指 $(Ku,v)=(u,Kv)$。对单射的正紧算子，规范本征函数因此构成整个空间的基；应用于适当平移后椭圆算子的逆，就得到离散本征展开。上述谱结论可参见文献\cite{hunter}附录4.B，本书不展开其完整证明。

这些工具分别提供不动点、弱解或谱分解。它们不会自动给出经典正则性；从空间中的存在结论回到逐点微分方程，还需正文相应章节中的正则性和迹论证。

